\subsubsection{Cost-effective Inference on RTX 4090} 
The NVIDIA GeForce RTX 4090 is a highly cost-effective GPU with 24G memory. However, through in-depth memory profiling and analysis, we identified memory insufficiency as the primary bottleneck to serve our model on it. To address this challenge, we developed a highly memory-efficient inference architecture and performs systematically optimizations. As a result, we successfully deployed and ran our 4.5B-parameter model on a single RTX 4090 GPU, and also support our largest 24B model on an 8$\times$RTX 4090 GPUs. In the following section, we briefly introduce the key optimization techniques.

\paragraph{Memory Optimization} 
To address the memory constraints of the RTX 4090, we used a variety of techniques to do systematically optimization:
 
\begin{itemize}
\item \emph{Quantization:} We adopt the same quantization strategy (WA8A SmoothQuant) as for streaming video generation.
\item \emph{KV-offload:} KV-offload is a technique that stores the KV cache in CPU memory by default and dynamically re-load it back to the GPU as needed. This approach significantly reduces peak GPU memory usage and is widely adopted in long-sequence processing for large language models (LLMs). In \magi, we also adopt this technique to effectively address memory constraints.

\item \emph{Hybrid Parallelism and Communication Optimization:} The above two optimizations only enable 4.5B model deployment on a single RTX 4090 GPU. However, the largest 24B model further requires multi-GPU parallelism. Unlike the streaming setting where we primarily adopt a Ulysses-based context-parallelism (CP) approach, deployment on RTX 4090 employs a hybrid strategy combining pipeline-parallelism (PP) and context-parallelism.

Specifically, pipeline-parallelism is used to partition model weights, while context-parallelism is used to partition activations. However, since the RTX 4090 utilizes PCIe for inter-GPU communication, both PP and CP suffer from communication-induced bubbles that degrade compute utilization, as measured by Model FLOPs Utilization (MFU). For PP, we mitigate this by interleaving tasks to overlap GPU idle. For context-parallelism, we initially adopted the Ulysses approach, but found that communication could not be fully overlapped with computation under PCIe constraints.

Therefore, we propose an enhancement to Ulysses called Context Shuffle Overlap (CSO)(Details in Sec.~\ref{sec:appendix_context_shuffle_overlap}), which scatters each chunk evenly across all GPUs, enabling finer-grained overlap between computation and communication than plain Ulysses. This strategy significantly improves MFU under the limited interconnect bandwidth of the RTX 4090. 
\end{itemize}

With the above optimizations, we constrained peak memory usage to 19.07GB for the 4.5B model on a single RTX 4090 GPU, and 19.29 GB for the 24B model on 8$\times$RTX 4090 GPUs. For the 24B model, the maximum MFU reached 66\%.




















